\section{Implementation and experimental evidence}
\label{sec:experiments}

\subsection{What was executed}

The archive contains two intentionally different implementations of the
mathematical operations. \code{prototype/search.py} uses SymPy 1.14.0 for exact
matrix solving, polynomial bases, division, and nullspaces.
\code{prototype/checkers.py} uses only the Python standard library, representing
polynomials as sparse exponent maps with \code{Fraction} coefficients. The
checker neither imports the search module nor calls a symbolic simplifier. It
can replay saved bundles with Python's site packages disabled.

The environment was Python 3.13.5 on x86-64 Linux. The deterministic corpus seed
is \code{2026091507}. Search time is measured once per case with
\code{perf\_counter\_ns}; checker time is the median of five warm invocations.
Process startup, imports, JSON parsing, and any Lean work are excluded from those
timings. This is a small deterministic research corpus, not a production
performance study.

\begin{table}[ht]
\centering\small
\begin{tabular}{@{}lrrrrrr@{}}
\toprule
Worker & Cases & Cert. & Cex. & Unknown & Search (ms) & Check (ms) \\
\midrule
Observable closure & 36 & 22 & 14 & 0 & 3.098 & 0.225 \\
Pullback ideals & 13 & 9 & 3 & 1 & 6.306 & 0.179 \\
Binomial sums & 10 & 7 & 0 & 3 & 47.632 & 0.313 \\
\bottomrule
\end{tabular}

\caption{Recorded outcomes. Timing columns are medians over the successful
instances in each family; the checker entry for each instance is itself a median
of five warm runs. Cex. means independently replayed counterexample.}
\label{tab:metrics}
\end{table}

Across 59 cases, 38 yield certificates, 17 yield concrete counterexamples, and
four retain \Unknown. The unknown cases are expected budget or template
controls, not failed cases reclassified as successes. The exact per-instance
results, options, and timings are in \code{results/summary.json},
\code{results/problems.json}, and \code{results/measurements.csv}.

\subsection{Corpus construction and what each part tests}

The 36 linear cases comprise 18 exact invariant systems in dimensions two
through seven, four product-equivalence systems related by invertible changes
of coordinates, seven delayed-error shift systems, six guaranteed nonzero
random controls, and a dedicated noncommuting two-action orientation test.
The positive invariant fixtures deliberately contain unobserved coordinates;
this tests demand-directed closure rather than forcing the basis to equal the
entire state space.

The 13 polynomial cases comprise the mutual-squares example, its budget control,
two mutants, four scaled product relations under square/cube actions, a Cassini
relation and its mutant, and three strengthening chains. The target-only
initial ideals are not supplied with the final auxiliary generators. In the
largest chain, five generators are discovered after four strict extensions.

The ten summation cases comprise seven successful binomial-moment recurrences,
two smaller-template versions that remain unknown, and a bounded order-one
attempt for the cubic-binomial sum. The numerical cross-checks evaluate each
accepted recurrence for $n=0,\ldots,30$. Those evaluations are independent
consistency tests; the universal argument is the polynomial certificate plus
Theorem~\ref{thm:sum}, not those 31 values.

No instance was imported from a benchmark and then represented as unseen.
These cases are intentionally constructed to distinguish exact scope, boundary
behavior, and implementation failure modes. They do not estimate coverage on
arbitrary Lean developments. Likewise, the linear and polynomial case counts
are not compared with the many different historical runs already stored in
Forge's repository.

\subsection{Selected measurements and certificate sizes}

The largest product-equivalence example has total state dimension ten and an
observable basis of size five. The search processes eleven queued rows. Its
recorded search time is $18.576$ ms and median checker time $1.785$ ms. The
mutual-squares ideal example takes $10.862$ ms to search and $0.214$ ms to replay;
its complete compact problem-and-certificate JSON bundle is 562 bytes.

The length-five polynomial strengthening chain uses five generators and 15
pullback reductions. Its search takes $55.934$ ms and replay $0.810$ ms. The
squared-binomial zeroth moment is found after five template trials with seven
unknown columns; its recorded search is $19.534$ ms and replay $0.261$ ms. The
binomial third moment requires total certificate degree six in the successful
run and takes $220.754$ ms to search.

All 38 successful compact JSON bundles total 27,814 bytes before ZIP
compression. The raw column named \code{certificate\_bytes} counts the entire
serialized bundle, including its problem and case name; it is not the size of a
Lean proof term or the certificate payload alone. Actual Lean reconstruction
and kernel-check costs remain unmeasured.

These small measurements support one modest conclusion: for this corpus,
finding an explanation is more expensive than replaying its finite identities,
and the saved explanations are compact. They do not support a general speedup
claim, an asymptotic scaling claim about SymPy's implementation, or a comparison
with \code{grind}.

\subsection{Independent replay, mutations, and schema controls}

The command
\begin{lstlisting}
python -S prototype/checkers.py results/certificates results/counterexamples
\end{lstlisting}
replays 55 bundles: 38 invariant or recurrence certificates and 17 concrete
counterexamples. The retained output reports zero failures and
\code{site\_packages\_enabled: false}. The four unknown outcomes have no proof or
counterexample certificate and are not included in this replay count.

For every successful certificate, the mutation harness adds one to each stored
matrix entry or nonzero polynomial coefficient, one position at a time. All
818 resulting mutations are rejected. Some mutations create a noncanonical
explicit zero monomial, while others preserve the encoding and break an
algebraic identity. This is a finite adversarial corpus, not a proof that every
possible forgery is rejected.

The 117 pytest items additionally cover target binding, altered initial states,
missing action coverage, duplicate JSON keys, floating-point tokens,
bool-as-integer confusion, unreduced rationals, duplicate or unordered
monomials, negative exponents, excessive declared degree, empty linear bases,
word direction, a nonlinear budget cutoff, zero binomial endpoints, and the
singular-recurrence nonuniqueness example. One item performs 60 independently
seeded sparse-arithmetic differential trials against SymPy. These are different
units: 117 test items are not 117 distinct mathematical theorems, and 818
mutations are not 818 new source goals.

\subsection{Failures retained rather than hidden}

Development exposed two substantive mistakes in the initial test expectations.
First, a randomly generated matrix system was labeled negative without ensuring
that its target ever became nonzero. The solver correctly certified it. The
fixture generator was repaired to force an observed nonzero update, and a
separate length-two case was added for word ordering. This was a fixture error,
not a failed mathematical certificate.

Second, the first moment-three summation expectation assumed that a degree-four
$r$ template would suffice. It did not. The smaller-template unknown remains in
the final corpus; the successful run explicitly permits degree six. The same
paired control was added for the squared-binomial second moment. Reporting only
the later successful run would obscure the actual search limitation.

Two test-development failures are also logged: a collection-time bracket error,
and an assertion that incorrectly required both recurrence coefficients to
vanish at the singular index. Only the leading coefficient needs to vanish to
lose the next-value constraint; the other coefficient can multiply a zero
current value. The corrected regression still checks two different sequences
with the same recurrence and first seed.

The retained diagnostic files explain these changes. They are not counted as
accepted runs, and they do not contain full historical source snapshots.
\code{docs/validation.md} states that distinction explicitly.

\subsection{Limitations of the executed artifact}

The Python checker has input limits: at most 16 scalar dimensions, bounded
integer encodings, at most 20,000 monomials per admitted polynomial, degree at
most 128, and a file-size ceiling. It checks intermediate polynomial growth as
well. These are pragmatic rejection conditions, not a formal resource theorem.
External process isolation remains necessary for adversarial use, particularly
because exact arithmetic and symbolic search can consume significant memory
before the next boundary check.

The supplied Lean code was not compiled, the source reifier was not built,
and none of the source-to-IR bridges was machine-verified. No claim is made that
117 Python tests remove the possibility of a checker bug. What is available
now is a reproducible mathematical implementation and a small, sharply specified
set of theorems that the Lean implementation should prove.
